\documentclass[a4paper,11pt,reqno]{amsart}

% =========================================================================
%  Instrumentación Electrónica — Sesión 1: Amplificadores operacionales
%  Apuntes de fundamentos y PreLab. Compilar con: pdflatex + biber + pdflatex ×2
%  (o latexmk -pdf main.tex).  Todo el preámbulo está en preambulo.tex.
% =========================================================================
% =========================================================================
%  Instrumentación Electrónica — Sesión 1: Amplificadores operacionales
%  Estilo heredado de los apuntes de Electromagnetismo I (Temas 0 y 1):
%  amsart, Palatino, versalitas, enunciados numerados por sección, paleta de
%  azules y grises sobrios. Ecuaciones numeradas a la derecha (opción reqno).
% =========================================================================

% --- IDIOMA Y CODIFICACIÓN ------------------------------------------------
\usepackage[utf8]{inputenc}
\usepackage[T1]{fontenc}
\usepackage[spanish,es-tabla,es-nodecimaldot]{babel}

% --- TIPOGRAFÍA (Palatino) --------------------------------------------------
\usepackage{mathpazo}
\linespread{1.05}
\emergencystretch=2em
\usepackage[final]{microtype}

% --- MATEMÁTICAS -------------------------------------------------------------
\usepackage{amsmath, amssymb, mathtools}
\usepackage{bm}

% --- UNIDADES (coma decimal, como en el resto de los apuntes) ---------------
\usepackage{siunitx}
\sisetup{output-decimal-marker={,}, group-separator={\,}, group-minimum-digits=5,
         exponent-product=\cdot, per-mode=symbol, range-phrase={--},
         list-final-separator={ y }, list-pair-separator={ y }}

% --- PÁGINA, COLORES Y GRÁFICOS --------------------------------------------
\usepackage[a4paper,margin=2.6cm]{geometry}
\usepackage{graphicx}
\usepackage{xcolor}
% Paleta base (idéntica a la de los apuntes de Electromagnetismo I)
\definecolor{azulprof}{HTML}{1F3A5F}    % azul oscuro
\definecolor{azulacero}{HTML}{3B6EA5}   % azul medio
\definecolor{azulclaro}{HTML}{8FB3D9}   % azul claro
\definecolor{azulfondo}{HTML}{F3F7FB}   % fondo tenue de definiciones
\definecolor{grispizarra}{HTML}{6B7280} % gris
\definecolor{grisclaro}{HTML}{D4D8DE}   % gris claro
\definecolor{acento}{HTML}{B8800A}      % acento: ámbar oscuro
\definecolor{rojoteo}{HTML}{B8800A}     % señales de salida, resultados
\definecolor{verdeteo}{HTML}{2F6F68}    % corrientes y sentidos de referencia
\definecolor{opfondo}{HTML}{E4EEF8}     % relleno del triángulo del AO

\usepackage{tikz}
\usetikzlibrary{babel, arrows.meta, calc, positioning, decorations.markings,
                patterns, intersections}
\usepackage[european, straightvoltages, RPvoltages]{circuitikz}
\usepackage{pgfplots}
\pgfplotsset{compat=1.17}
\usepgfplotslibrary{fillbetween}
\usepackage{booktabs}
\usepackage{needspace}
\usepackage{array}
\usepackage{multirow}
\usepackage{enumitem}
\setlist[itemize]{leftmargin=1.6em, itemsep=2pt, topsep=4pt}
\setlist[enumerate]{leftmargin=1.8em, itemsep=2pt, topsep=4pt}

% --- ENLACES Y REFERENCIAS CRUZADAS ----------------------------------------
\usepackage[colorlinks=true, linkcolor=azulacero, citecolor=azulacero,
            urlcolor=azulacero, plainpages=false, pdfusetitle=false,
            pdftitle={Sesión 1. Amplificadores operacionales},
            pdfauthor={Daniel Martín Guerrero}]{hyperref}
\usepackage[spanish,nameinlink]{cleveref}

% --- BIBLIOGRAFÍA -----------------------------------------------------------
\usepackage{csquotes}
\usepackage[backend=biber, style=numeric, sorting=none, giveninits=true,
            maxbibnames=6, minbibnames=3, doi=true, url=true,
            isbn=false]{biblatex}
\addbibresource{referencias.bib}
\DeclareFieldFormat[article,techreport,unpublished,book,manual]{title}{#1\isdot}
\DeclareFieldFormat{url}{\url{#1}}
\renewcommand*{\bibfont}{\small}
\setlength{\bibitemsep}{3pt}

% --- CABECERAS ---------------------------------------------------------------
\pagestyle{headings}
\makeatletter
\def\ps@headings{%
  \let\@mkboth\markboth
  \def\@evenhead{\small\scshape\normalfont\thepage\hfil Instrumentación Electrónica}%
  \def\@oddhead{\small\scshape\normalfont Sesión 1. Amplificadores operacionales\hfil\thepage}%
  \let\@evenfoot\@empty \let\@oddfoot\@empty}
\makeatother
\numberwithin{equation}{section}
% Subsecciones como títulos en bloque (no «run-in»), con aire antes y después
\usepackage{titlesec}
\titleformat{\subsection}[block]{\normalfont\bfseries\color{azulprof}}{\thesubsection.}{0.6em}{}
\titlespacing*{\subsection}{0pt}{3ex plus 1ex minus .5ex}{1.6ex plus .4ex}
\setcounter{tocdepth}{2}
\makeatletter
\def\@@and{y}
\renewcommand{\keywordsname}{Palabras clave}
% Índice: número de sección en una caja de anchura fija (evita el solape de amsart)
\renewcommand{\tocsection}[3]{\indentlabel{\@ifnotempty{#2}{\makebox[1.6em][l]{#2.}}}#3}
\renewcommand{\tocsubsection}[3]{\indentlabel{\@ifnotempty{#2}{\makebox[2.4em][l]{#2.}}}#3}
\def\l@section{\@tocline{1}{2pt}{0pt}{1.6em}{}}
\def\l@subsection{\@tocline{2}{0pt}{1.6em}{2.4em}{}}
\def\l@part{\@tocline{-1}{8pt plus 2pt}{0pt}{0pt}{\bfseries\color{azulprof}}}
\makeatother

% =========================================================================
%               NOTACIÓN Y MACROS DE USO FRECUENTE
% =========================================================================
\newcommand{\dif}{\mathrm{d}}
\newcommand{\eul}{\mathrm{e}}
\newcommand{\jj}{\mathrm{j}}                         % unidad imaginaria (j)
\newcommand{\R}{\mathbb{R}}
\newcommand{\C}{\mathbb{C}}
\newcommand{\abs}[1]{\lvert #1\rvert}
\newcommand{\defeq}{\mathrel{\overset{\mathrm{def}}{=}}}
\newcommand{\Lt}{\mathcal{L}}                        % transformada de Laplace
\DeclareMathOperator{\sen}{sen}
\DeclareMathOperator{\sgn}{sgn}
\DeclareMathOperator{\arctg}{arctg}
\DeclareMathOperator{\Real}{Re}
\renewcommand{\sin}{\sen}                            % «sen» en todo el texto
% Magnitudes de circuito
\newcommand{\vin}{v_{\mathrm{in}}}
\newcommand{\vout}{v_{\mathrm{out}}}
\newcommand{\Vin}{V_{\mathrm{in}}}
\newcommand{\Vout}{V_{\mathrm{out}}}
\newcommand{\Vsat}{V_{\mathrm{sat}}}
\newcommand{\Vcc}{V_{\mathrm{CC}}}
\newcommand{\Vee}{V_{\mathrm{EE}}}
\newcommand{\Rin}{R_{\mathrm{in}}}
\newcommand{\Rout}{R_{\mathrm{out}}}
\newcommand{\Zin}{Z_{\mathrm{in}}}
\newcommand{\Zout}{Z_{\mathrm{out}}}
\newcommand{\Av}{A_{V}}
\newcommand{\vp}{v_{+}}
\newcommand{\vm}{v_{-}}
\newcommand{\dB}{\,\mathrm{dB}}
\newcommand{\pend}[1]{\textcolor{grispizarra}{\colorbox{grisclaro}{\strut\,#1\,}}}

% =========================================================================
%        ENTORNOS FORMALES (enunciados numerados, estilo artículo)
% =========================================================================
\usepackage{amsthm}
\usepackage[most]{tcolorbox}

\newtheoremstyle{cienplano}{10pt}{10pt}{\itshape}{}%
  {\bfseries\color{azulprof}}{.}{0.6em}%
  {\thmname{#1}\thmnumber{ #2}\thmnote{\normalfont\upshape\ (#3)}}
\newtheoremstyle{ciendef}{10pt}{10pt}{\normalfont}{}%
  {\bfseries\color{azulprof}}{.}{0.6em}%
  {\thmname{#1}\thmnumber{ #2}\thmnote{\normalfont\upshape\ (#3)}}
\newtheoremstyle{cienobs}{8pt}{8pt}{\normalfont}{}%
  {\itshape\color{azulprof}}{.}{0.6em}%
  {\thmname{#1}\thmnumber{ #2}\thmnote{\normalfont\upshape\ (#3)}}

\theoremstyle{cienplano}
\newtheorem{proposicion}{Proposición}[section]
\newtheorem{teorema}[proposicion]{Teorema}
\newtheorem{corolario}[proposicion]{Corolario}
\theoremstyle{ciendef}
\newtheorem{ejemplo}[proposicion]{Ejemplo}
\newtheorem{ejercicio}[proposicion]{PreLab}
\newtheorem{hipotesis}[proposicion]{Hipótesis}
% Definiciones: único entorno enmarcado (caja fina, fondo tenue, título en el borde)
\newcounter{definicion}
\renewcommand{\thedefinicion}{\theproposicion}
\newtcolorbox{definicionbox}[1]{%
  enhanced, breakable, sharp corners,
  colback=azulfondo, colframe=azulclaro, boxrule=0.4pt,
  left=9pt, right=9pt, top=5pt, bottom=5pt,
  before skip=12pt, after skip=12pt,
  attach boxed title to top left={yshift*=-\tcboxedtitleheight/2, xshift=8pt},
  boxed title style={colback=white, colframe=white, sharp corners, boxrule=0pt,
                     left=3pt, right=3pt, top=1pt, bottom=1pt},
  fonttitle=\bfseries\color{azulprof}, coltitle=azulprof,
  title={Definición~\thedefinicion\ifx\relax#1\relax\else\ (#1)\fi}}
\newenvironment{definicion}[1][]
  {\par\stepcounter{proposicion}\refstepcounter{definicion}\begin{definicionbox}{#1}}
  {\end{definicionbox}}

% Enunciado de cada recuadro del PreLab: misma caja, con filete ámbar (los
% recuadros amarillos del guion), para distinguir lo pedido de lo resuelto.
\newcounter{prelab}
\newtcolorbox{prelabbox}[1]{%
  enhanced, breakable, sharp corners,
  colback=white, colframe=acento!70, boxrule=0.5pt, leftrule=2.2pt,
  left=9pt, right=9pt, top=5pt, bottom=5pt,
  before skip=12pt, after skip=10pt,
  attach boxed title to top left={yshift*=-\tcboxedtitleheight/2, xshift=8pt},
  boxed title style={colback=white, colframe=white, sharp corners, boxrule=0pt,
                     left=3pt, right=3pt, top=1pt, bottom=1pt},
  fonttitle=\bfseries\color{acento!85!black}, coltitle=acento!85!black,
  title={PreLab~\theprelab\ifx\relax#1\relax\else\ (#1)\fi}}
\newenvironment{prelab}[1][]
  {\par\refstepcounter{prelab}\begin{prelabbox}{#1}}
  {\end{prelabbox}}

\theoremstyle{cienobs}
\newtheorem{observacion}[proposicion]{Observación}

% Resolución: misma maquinaria que «proof», con otro rótulo
\renewcommand{\qedsymbol}{\ensuremath{\blacksquare}}
\newenvironment{resolucion}[1][Resolución]
  {\begin{proof}[\itshape\color{azulprof}#1]}{\end{proof}}

% Resultado final enmarcado (como \boxed, pero con el filete de la paleta)
\newcommand{\resultado}[1]{\fcolorbox{azulacero}{white}{\ensuremath{\displaystyle #1}}}

\crefname{proposicion}{proposición}{proposiciones}
\Crefname{proposicion}{Proposición}{Proposiciones}
\crefname{teorema}{teorema}{teoremas}
\Crefname{teorema}{Teorema}{Teoremas}
\crefname{corolario}{corolario}{corolarios}
\Crefname{corolario}{Corolario}{Corolarios}
\crefname{definicion}{definición}{definiciones}
\Crefname{definicion}{Definición}{Definiciones}
\crefname{ejemplo}{ejemplo}{ejemplos}
\Crefname{ejemplo}{Ejemplo}{Ejemplos}
\crefname{hipotesis}{hipótesis}{hipótesis}
\Crefname{hipotesis}{Hipótesis}{Hipótesis}
\crefname{prelab}{PreLab}{PreLabs}
\Crefname{prelab}{PreLab}{PreLabs}
\crefname{observacion}{observación}{observaciones}
\Crefname{observacion}{Observación}{Observaciones}
\crefname{section}{sección}{secciones}
\Crefname{section}{Sección}{Secciones}
\crefname{subsection}{sección}{secciones}
\Crefname{subsection}{Sección}{Secciones}
\crefname{table}{tabla}{tablas}
\Crefname{table}{Tabla}{Tablas}
\crefname{figure}{figura}{figuras}
\Crefname{figure}{Figura}{Figuras}
\crefname{equation}{ecuación}{ecuaciones}
\Crefname{equation}{Ecuación}{Ecuaciones}
\crefformat{equation}{ecuación~(#2#1#3)}
\Crefformat{equation}{Ecuación~(#2#1#3)}
\crefrangeformat{equation}{ecuaciones~(#3#1#4)--(#5#2#6)}
\crefmultiformat{equation}{ecuaciones~(#2#1#3)}{ y~(#2#1#3)}{, (#2#1#3)}{ y~(#2#1#3)}

\usepackage[font=small, labelfont={sc,color=azulprof}, labelsep=period,
            margin=0.6cm, justification=justified]{caption}
\usepackage{subcaption}

% Referencia a un panel de figura: «figura 5\,(b)», enlazada en bloque
\newcommand{\figpanel}[2]{\hyperref[#1]{figura~\ref*{#1}\,\textup{(#2)}}}
\newcommand{\Figpanel}[2]{\hyperref[#1]{Figura~\ref*{#1}\,\textup{(#2)}}}

% =========================================================================
%        ESTILOS GRÁFICOS COMUNES (TikZ / circuitikz / pgfplots)
% =========================================================================
% =========================================================================
%  Estilos gráficos comunes: TikZ, circuitikz y pgfplots
% =========================================================================

% --- Estilos TikZ generales ------------------------------------------------
\tikzset{
  >=Stealth,
  eje/.style={->, line width=0.6pt, grispizarra},
  fino/.style={line width=0.4pt, grispizarra},
  punteado/.style={densely dotted, line width=0.5pt, grispizarra},
  etiq/.style={font=\small, text=black!85},
  nota/.style={font=\scriptsize, text=grispizarra, align=center},
  % corriente dibujada junto a un hilo (flecha verde con rótulo)
  corr/.style={-{Stealth[length=4.5pt,width=4pt]}, line width=0.8pt, verdeteo},
  % nudo con tensión anotada
  tens/.style={font=\small, text=azulprof},
}

% --- circuitikz: trazo y colores sobrios -----------------------------------
\ctikzset{
  bipoles/thickness=1,
  bipoles/length=1.15cm,
  tripoles/op amp/width=1.5,
  tripoles/op amp/height=1.35,
  tripoles/op amp/font=\small,
  current/distance=0.45,
  label/align=smart,
}
\tikzset{
  circuito/.style={line width=0.55pt, color=black!85, font=\small},
  ao/.style={fill=opfondo},
}

% --- Estilo base de ejes pgfplots ------------------------------------------
\pgfplotsset{
  tema/.style={
    tick label style={font=\scriptsize}, label style={font=\small},
    axis lines=left, axis line style={grispizarra},
    grid=major, grid style={grisclaro, line width=0.4pt},
    every tick/.style={grispizarra},
    legend style={font=\scriptsize, draw=none, fill=white, fill opacity=0.85,
                  text opacity=1},
    legend cell align=left,
    /pgf/number format/use comma,
  },
  % ejes centrados (formas de onda y características de transferencia)
  temacentro/.style={
    tick label style={font=\scriptsize}, label style={font=\small},
    axis lines=middle, axis line style={grispizarra, -{Stealth[length=4pt]}},
    every tick/.style={grispizarra},
    legend style={font=\scriptsize, draw=none, fill=white, fill opacity=0.85,
                  text opacity=1},
    legend cell align=left,
    /pgf/number format/use comma,
  },
  senal entrada/.style={azulacero, line width=1.0pt},
  senal salida/.style={rojoteo, line width=1.1pt},
  asintota/.style={grispizarra, densely dashed, line width=0.6pt},
}



% =========================================================================
%                     VARIABLES PARA LA PORTADA
% =========================================================================
\newcommand{\Universidad}{Universidad de Salamanca}
\newcommand{\Facultad}{Facultad de Ciencias}
\newcommand{\Titulacion}{Grado en Física}
\newcommand{\Asignatura}{Instrumentación Electrónica}
\newcommand{\NumTema}{Laboratorio de circuitos, sesión 1}
\newcommand{\TituloTema}{Amplificadores\\operacionales}
\newcommand{\tipotrabajo}{Apuntes de fundamentos y PreLab}
\newcommand{\miNombre}{Daniel Martín Guerrero}
\newcommand{\miRotuloProfesor}{Profesores}
\newcommand{\miProfesor}{Yahya Moubarak Meziani\\María Susana Pérez Santos\\Jesús Enrique Velázquez Pérez}
\newcommand{\miDetalleProfesor}{Área de Electrónica}
\newcommand{\SubtituloTema}{Amplificadores inversor y no inversor,\\amplificador inversor generalizado y filtro paso bajo de primer orden}
\newcommand{\miUbicacion}{Salamanca, curso 2026/2027}

% =========================================================================
%                     DATOS DEL ARTÍCULO (cabecera)
% =========================================================================
\title{Amplificadores operacionales: fundamentos y PreLab}
\author{Daniel Martín Guerrero}
\address{Grado en Física, Facultad de Ciencias, Universidad de Salamanca,
         Salamanca, España}
\date{}

% Separador de las dos partes del documento (entra también en el índice)
\newcommand{\parte}[2]{%
  \clearpage
  \phantomsection
  \addcontentsline{toc}{part}{Parte #1. #2}%
  \begin{center}
    {\color{azulacero}\rule{\textwidth}{0.5pt}}\\[6pt]
    {\color{grispizarra}\scshape\small\textls[120]{Parte #1}}\\[2pt]
    {\color{azulprof}\scshape\large\textls[80]{#2}}\\[2pt]
    {\color{azulacero}\rule{\textwidth}{0.5pt}}
  \end{center}
  \vspace{4pt}}

% =========================================================================
%                                DOCUMENTO
% =========================================================================
\begin{document}
% --- PORTADA ----------------------------------------------------------------
\pagenumbering{alph}
%%%%%%%%%%%%%%%%%%%%%%%%%%%%%%%%%%%%%%%%%%%%%%%%%%%%%%%%%%%%%%%%%%%%%%%%
% Portada formal (estilo tesis): logotipos, identificación institucional,
% título y datos de autoría. Banda vertical a la derecha, como en el
% «Cálculo integral multivariable» de L. Navas. Requiere en main.tex las
% macros \Universidad, \Facultad, \Titulacion, \Asignatura, \NumTema,
% \TituloTema, \SubtituloTema, \tipotrabajo, \miNombre, \miProfesor,
% \miUbicacion, \miRotuloProfesor, \miDetalleProfesor.
%%%%%%%%%%%%%%%%%%%%%%%%%%%%%%%%%%%%%%%%%%%%%%%%%%%%%%%%%%%%%%%%%%%%%%%%

\begin{titlepage}
\thispagestyle{empty}
\pagecolor{white}

\begin{tikzpicture}[remember picture, overlay]
  \coordinate (O) at (current page.south west);

  % Velo muy tenue en la mitad superior (cabecera institucional)
  \shade[top color=azulfondo, bottom color=white]
    ($(O)+(0,18.5cm)$) rectangle ($(O)+(21cm,29.7cm)$);

  % Banda vertical azul, a la derecha del bloque de texto
  \fill[azulacero] ($(O)+(19.0cm,2.6cm)$) rectangle ($(O)+(19.14cm,27.4cm)$);

  % Logotipos
  \node[anchor=north west, inner sep=0] at ($(O)+(2.4cm,27.7cm)$)
    {\includegraphics[height=1.55cm]{logo_usal.png}};
  \node[anchor=north east, inner sep=0] at ($(O)+(17.4cm,27.7cm)$)
    {\includegraphics[height=1.55cm]{logo_facultad.png}};
  \draw[grispizarra!60, line width=0.4pt] ($(O)+(2.4cm,25.7cm)$) -- ($(O)+(17.4cm,25.7cm)$);

  % Identificación institucional
  \node[anchor=north, align=center, inner sep=0, text width=15cm] at ($(O)+(9.9cm,25.2cm)$)
    {{\color{azulprof}\scshape\Large\textls[120]{\Universidad}}\\[5pt]
     {\color{grispizarra}\scshape\large\textls[80]{\Facultad}}\\[3pt]
     {\color{grispizarra}\scshape\large\textls[80]{\Titulacion}}};

  % Asignatura y título
  \node[anchor=north, align=center, inner sep=0, text width=15cm] at ($(O)+(9.9cm,19.2cm)$)
    {{\color{azulacero}\scshape\Large\textls[160]{\Asignatura}}\\[4pt]
     {\color{grispizarra}\scshape\large\textls[100]{\NumTema}}\\[22pt]
     {\color{azulprof}\fontsize{34}{40}\selectfont\bfseries\TituloTema}\\[14pt]
     {\color{azulacero}\rule{3.2cm}{0.8pt}}\\[14pt]
     {\color{black!80}\itshape\large\SubtituloTema}};

  % Epígrafe: los tres resultados del PreLab
  \node[anchor=north, align=center, inner sep=0, text width=15cm] at ($(O)+(9.9cm,11.0cm)$)
    {\color{grispizarra}\small
     $\displaystyle A_{V}=-\frac{R_{2}}{R_{1}}\qquad
      A_{V}=1+\frac{R_{2}}{R_{1}}\qquad
      H(s)=-\frac{Z_{2}(s)}{Z_{1}(s)}$};

  % Autoría y dirección
  \draw[grispizarra!60, line width=0.4pt] ($(O)+(2.4cm,8.6cm)$) -- ($(O)+(17.4cm,8.6cm)$);
  \node[anchor=north west, inner sep=0, text width=7cm, align=left] at ($(O)+(2.4cm,8.2cm)$)
    {{\scshape\small\color{azulacero}\textls[80]{Autor}}\\[4pt]
     {\large\color{black!85}\miNombre}\\[2pt]
     {\small\color{grispizarra}Grado en Física}};
  \node[anchor=north west, inner sep=0, text width=7cm, align=left] at ($(O)+(10.4cm,8.2cm)$)
    {{\scshape\small\color{azulacero}\textls[80]{\miRotuloProfesor}}\\[4pt]
     {\color{black!85}\miProfesor}\\[3pt]
     {\small\color{grispizarra}\miDetalleProfesor}};

  % Tipo de documento, lugar y curso
  \node[anchor=south, align=center, inner sep=0, text width=15cm] at ($(O)+(9.9cm,2.9cm)$)
    {{\color{azulprof}\scshape\large\textls[100]{\tipotrabajo}}\\[4pt]
     {\color{grispizarra}\scshape\small\textls[80]{\miUbicacion}}};
\end{tikzpicture}

\end{titlepage}

\clearpage
\pagenumbering{arabic}
\setcounter{page}{1}

\begin{abstract}
Reunimos, de forma sintética, los conceptos de electrónica analógica que exige
la primera sesión del laboratorio de circuitos: señales y sus valores
característicos, el amplificador como cuadripolo (impedancias de entrada y de
salida, ganancias y decibelios), la impedancia compleja, la transformada de
Laplace y la función de transferencia, el diagrama de Bode de una red de
constante temporal única y, por último, el amplificador operacional con sus
modelos ideal y real. Con estas herramientas resolvemos, paso a paso y con
cada hipótesis justificada, los recuadros del PreLab: diseño de los
amplificadores inversor ($A_V=-10$) y no inversor ($A_V=+11$), y obtención de la
función de transferencia del amplificador inversor generalizado, que con
$Z_1=R$ y $Z_2=R_F\parallel C$ se comporta como un filtro paso bajo de primer
orden.
\end{abstract}

\maketitle
\tableofcontents

% =========================================================================
\parte{I}{Fundamentos}
% =========================================================================
% =========================================================================
\section{Introducción: qué estudiamos en la sesión}\label{sec:introduccion}
% =========================================================================
Un \emph{amplificador operacional} (AO) es un amplificador de tensión diferencial con una
ganancia tan grande que, por sí solo, apenas es útil: cualquier diferencia de milivoltios
entre sus entradas lo lleva a saturación. Su interés está en la \emph{realimentación
negativa}: si se devuelve una fracción de la salida a la entrada inversora mediante una red
pasiva, el comportamiento del conjunto deja de depender del AO y pasa a depender solo de esa
red. La sesión~1 del laboratorio~\cite{GUION-S1} explota esta idea en tres circuitos
(\cref{fig:esquema}):
\begin{enumerate}
  \item un \emph{amplificador inversor} de ganancia $A_V=-10$;
  \item un \emph{amplificador no inversor} de ganancia $A_V=+11$;
  \item un \emph{amplificador inversor generalizado}, en el que las resistencias se sustituyen
        por impedancias. Si la red de realimentación contiene un condensador, la ganancia
        depende de la frecuencia y el circuito es un \emph{filtro}; en el dominio del tiempo
        integra o deriva la señal, y la transformada de Laplace enlaza ambas descripciones.
\end{enumerate}
La parte~I reúne las herramientas imprescindibles, con la notación de la teoría de la
asignatura~\cite{TEORIA-EA,TEORIA-ADICIONAL,TEORIA-LAPLACE} y del texto de referencia
de Sedra y Smith~\cite{SEDRA}; la parte~II resuelve los recuadros del PreLab.

% Figura: estructura lógica de la sesión.
\begin{figure}[ht]
\centering
\begin{tikzpicture}[
    caja/.style={draw=azulclaro, fill=azulfondo, line width=0.5pt, align=center,
                 inner sep=5pt, font=\small, minimum height=1.0cm},
    final/.style={draw=azulacero, fill=white, line width=0.8pt, align=center,
                  inner sep=5pt, font=\small, minimum height=1.0cm},
    fl/.style={->, line width=0.7pt, azulacero, shorten >=2pt, shorten <=2pt},
    nota/.style={font=\scriptsize, text=grispizarra, align=center}]
  \node[caja] (M) at (0,0) {Modelo ideal del AO\\[-1pt]{\footnotesize $i_\pm=0,\ v_+=v_-$}};
  \node[caja] (L) at (0,-1.9) {Impedancias\\[-1pt]{\footnotesize $Z_R=R,\ Z_C=1/(sC)$}};
  \node[caja] (G) at (4.4,-0.95) {Inversor generalizado\\[-1pt]{\footnotesize $H(s)=-Z_2/Z_1$}};
  \node[final] (I) at (9.3,-0.95) {Inversor\\[-1pt]{\footnotesize $A_V=-R_2/R_1$}};
  \node[final] (N) at (9.3,0.85) {No inversor\\[-1pt]{\footnotesize $A_V=1+R_2/R_1$}};
  \node[final] (F) at (9.3,-2.65) {Filtro paso bajo\\[-1pt]{\footnotesize $Z_2=R_F\parallel C$}};
  \draw[fl] (M.east) -- (G.north west);
  \draw[fl] (L.east) -- (G.south west);
  \draw[fl] (G.east) -- (I.west);
  \draw[fl] (G.east) -- (F.west);
  \draw[fl] (M.east) ++(0,0.12) -- ($(N.west)+(0,0)$);
  \node[nota] at (6.95,-0.72) {$Z_k=R_k$};
  \node[nota] at (6.75,-2.25) {$Z_1=R$};
\end{tikzpicture}
\caption{Estructura lógica de la sesión. El modelo ideal del AO, junto con el concepto de
impedancia, da la ganancia del amplificador inversor generalizado; el inversor y el filtro
paso bajo son casos particulares. El no inversor se resuelve con las mismas reglas.}
\label{fig:esquema}
\end{figure}


% =========================================================================
\section{Señales y amplificadores}\label{sec:amplificadores}
% =========================================================================
% -------------------------------------------------------------------------
\subsection{Valores característicos de una señal}\label{sec:senales}
% -------------------------------------------------------------------------
En el laboratorio una señal es una tensión $y(t)$ que el osciloscopio representa frente al tiempo.
Para compararla con los cálculos se resume en unos pocos números~\cite{TEORIA-EA}.

\begin{definicion}[Valores de una señal periódica]\label{def:valores}
Sea $y(t)$ periódica de periodo $T=1/f$. Se definen el valor de pico, el valor pico a pico, el
valor medio y el valor eficaz (\emph{root mean square}):
\begin{gather}
  y_{p}=\max_{t}\abs{y(t)},\qquad
  y_{pp}=\max_{t}y(t)-\min_{t}y(t),\label{eq:pico}\\
  y_{m}=\frac1T\int_{0}^{T}y(t)\,\dif t,\qquad
  y_{\mathrm{rms}}=\biggl(\frac1T\int_{0}^{T}y^{2}(t)\,\dif t\biggr)^{1/2}.\label{eq:rms}
\end{gather}
\end{definicion}

Para una señal simétrica de media nula, $y_{pp}=2y_{p}$, y el valor eficaz vale $y_p/\sqrt2$
(sinusoide), $y_p$ (onda cuadrada) o $y_p/\sqrt3$ (onda triangular). El osciloscopio mide con
comodidad $y_{pp}$; por eso, en lo que sigue, las ganancias se compararán como cocientes de
amplitudes o de tensiones pico a pico, que coinciden para señales de la misma forma.

% -------------------------------------------------------------------------
\subsection{El amplificador como cuadripolo}\label{sec:cuadripolo}
% -------------------------------------------------------------------------
Cualquier amplificador lineal, visto desde sus dos puertos (entrada y salida, con un terminal
común a tierra), se describe con el modelo de la \cref{fig:cuadripolo}: la entrada es una
impedancia $\Zin$ y la salida, un equivalente de Thévenin formado por una fuente controlada
$A_{vo}v_i$ y una impedancia $\Zout$~\cite{TEORIA-EA,TEORIA-ADICIONAL,SEDRA}.

% Figura: modelo de cuadripolo de un amplificador de tensión, con fuente y carga.
\begin{figure}[ht]
\centering
\begin{circuitikz}[circuito, scale=0.95, transform shape]
  % Recuadro del amplificador
  \fill[azulfondo] (2.8,-0.45) rectangle (7.3,2.65);
  \draw[azulclaro, line width=0.5pt, densely dashed] (2.8,-0.45) rectangle (7.3,2.65);
  \node[nota, anchor=south] at (5.05,2.65) {amplificador};
  % Fuente y resistencia interna
  \draw (0,0) to[sV, l=$v_s$] (0,2) to[R, l=$R_s$] (2.2,2)
        to[short, o-, i=$i_i$] (3.5,2) to[R, l=$\Rin$] (3.5,0) -- (2.2,0) to[short, -o] (2.2,0) -- (0,0);
  \draw (2.2,2) to[open, v^=$v_i$] (2.2,0);
  % Salida: equivalente de Thévenin
  \draw (3.5,0) -- (5.3,0) to[cV, l_=$A_{vo}v_i$] (5.3,2) to[R, l=$\Rout$] (6.9,2)
        to[short, -o, i=$i_o$] (7.9,2) -- (9.1,2) to[R, l=$R_L$] (9.1,0) -- (7.9,0)
        to[short, o-] (7.9,0) -- (5.3,0);
  \draw (7.9,2) to[open, v^=$v_o$] (7.9,0);
  \draw (4.4,0) node[ground] {};
  % Impedancias vistas desde cada puerto
  \draw[->, azulacero, line width=0.6pt] (1.6,2.85) -- node[above, font=\scriptsize, text=azulacero] {$\Zin$} (2.6,2.85);
  \draw[->, azulacero, line width=0.6pt] (8.5,2.85) -- node[above, font=\scriptsize, text=azulacero] {$\Zout$} (7.5,2.85);
\end{circuitikz}
\caption{Modelo de un amplificador de tensión como cuadripolo~\cite{TEORIA-EA,SEDRA}. La fuente
$v_s$ tiene resistencia interna $R_s$ (en el generador del laboratorio, $50\,\Omega$) y el
amplificador alimenta una carga $R_L$. $\Zin$ es la impedancia que ve la fuente y $\Zout$ la que ve
la carga.}
\label{fig:cuadripolo}
\end{figure}


\begin{definicion}[Parámetros de un amplificador]\label{def:parametros}
Con las magnitudes de la \cref{fig:cuadripolo},
\begin{equation}\label{eq:parametros}
  \Zin=\frac{v_i}{i_i},\qquad
  A_V=\frac{v_o}{v_i},\qquad
  A_I=\frac{i_o}{i_i},\qquad
  A_W=\frac{v_o\,i_o}{v_i\,i_i}=A_V\,A_I ,
\end{equation}
y $\Zout$ es la impedancia de Thévenin vista desde los terminales de salida con la fuente de
entrada anulada ($v_s=0$).
$A_{vo}$ es la ganancia en tensión \emph{en circuito abierto} ($R_L\to\infty$).
\end{definicion}

\begin{proposicion}[Efecto de carga]\label{prop:carga}
Si el amplificador se excita con una fuente de resistencia interna $R_s$ y se carga con $R_L$, la
ganancia efectiva desde la fuente es
\begin{equation}\label{eq:carga}
  \frac{v_o}{v_s}=\underbrace{\frac{\Rin}{\Rin+R_s}}_{\text{divisor de entrada}}\;A_{vo}\;
                  \underbrace{\frac{R_L}{R_L+\Rout}}_{\text{divisor de salida}} .
\end{equation}
\end{proposicion}
\begin{proof}
En la entrada, $R_s$ y $\Rin$ forman un divisor de tensión: $v_i=v_s\Rin/(\Rin+R_s)$. En la salida, la
fuente $A_{vo}v_i$ ve en serie $\Rout$ y $R_L$: $v_o=A_{vo}v_i\,R_L/(R_L+\Rout)$. Se multiplican ambas.
\end{proof}

\begin{corolario}[Amplificador de tensión ideal]\label{cor:ideal}
La ganancia no depende de la fuente ni de la carga si y solo si $\Rin\gg R_s$ y $\Rout\ll R_L$; en el
límite ideal, $\Rin\to\infty$ y $\Rout\to0$, y entonces $v_o/v_s=A_{vo}$.
\end{corolario}

Por eso hay que conocer \emph{las dos} impedancias además de una ganancia: la ganancia en potencia
$A_W=A_V^{2}R_{\mathrm{in}}/R_L$ (con $i_o=v_o/R_L$, $i_i=v_i/\Rin$) depende de su cociente.

\begin{definicion}[Ganancias en decibelios]\label{def:dB}
\begin{equation}\label{eq:dB}
  A_V(\mathrm{dB})=20\log_{10}\abs{A_V},\qquad
  A_I(\mathrm{dB})=20\log_{10}\abs{A_I},\qquad
  A_W(\mathrm{dB})=10\log_{10}A_W .
\end{equation}
\end{definicion}
El factor $20$ aparece porque la potencia es cuadrática en la tensión: así, $A_W(\mathrm{dB})=A_V(\mathrm{dB})$
cuando $\Rin=R_L$. Valores útiles: $\abs{A_V}=10\mapsto20\dB$, $100\mapsto40\dB$, $1\mapsto0\dB$ y
$1/\sqrt2\mapsto-3{,}01\dB$. Las ganancias pueden ser negativas (inversión de fase), de ahí el módulo; y,
salvo a baja frecuencia, son números complejos, de modo que hay que dar también su fase.

% =========================================================================
\section{Impedancia, transformada de Laplace y respuesta en frecuencia}\label{sec:laplace}
% =========================================================================
El circuito~3 del guion contiene un condensador, cuya relación tensión--corriente es
diferencial. Esta sección reúne las tres ideas que permiten tratarlo con la misma
álgebra que una resistencia: la impedancia compleja, la transformada de Laplace y el
diagrama de Bode~\cite{TEORIA-LAPLACE,TEORIA-ADICIONAL}.

% -------------------------------------------------------------------------
\subsection{Régimen sinusoidal e impedancia compleja}\label{sec:impedancia}
% -------------------------------------------------------------------------
En un circuito lineal excitado por una sinusoide de pulsación $\omega=2\pi f$, todas las
tensiones y corrientes son, en régimen permanente, sinusoides de la misma pulsación. Cada una
queda determinada por su amplitud y su fase, que se agrupan en un número complejo, el \emph{fasor}:
\begin{equation}\label{eq:fasor}
  v(t)=V_p\cos(\omega t+\varphi)=\Real\bigl\{\mathbf{V}\,\eul^{\jj\omega t}\bigr\},\qquad
  \mathbf{V}=V_p\,\eul^{\jj\varphi}\in\C .
\end{equation}
Como $\frac{\dif}{\dif t}\eul^{\jj\omega t}=\jj\omega\,\eul^{\jj\omega t}$, \emph{derivar en el tiempo equivale a
multiplicar el fasor por $\jj\omega$}.

\begin{definicion}[Impedancia]\label{def:impedancia}
La impedancia de un dipolo lineal es el cociente de los fasores de su tensión y su corriente,
$Z(\jj\omega)=\mathbf{V}/\mathbf{I}$. Su módulo es el cociente de amplitudes y su argumento, el
desfase de la tensión respecto de la corriente.
\end{definicion}

\begin{proposicion}[Impedancias elementales]\label{prop:impedancias}
\begin{equation}\label{eq:impedancias}
  Z_R=R,\qquad Z_C=\frac{1}{\jj\omega C},\qquad Z_L=\jj\omega L .
\end{equation}
\end{proposicion}
\begin{proof}
Ley de Ohm, $v=Ri\Rightarrow\mathbf{V}=R\mathbf{I}$. En el condensador, $i=C\,\dif v/\dif t\Rightarrow
\mathbf{I}=\jj\omega C\,\mathbf{V}$. En la bobina, $v=L\,\dif i/\dif t\Rightarrow\mathbf{V}=\jj\omega L\,\mathbf{I}$.
\end{proof}

Las leyes de Kirchhoff son lineales, de modo que se cumplen también para los fasores; en
consecuencia, las impedancias se asocian como resistencias: en serie se suman,
$Z=Z_1+Z_2$, y en paralelo se suman sus inversas, $Z=Z_1Z_2/(Z_1+Z_2)$. En particular,
\begin{equation}\label{eq:RC-paralelo}
  Z_{R\parallel C}=\frac{R\cdot\frac{1}{\jj\omega C}}{R+\frac{1}{\jj\omega C}}=\frac{R}{1+\jj\omega RC} .
\end{equation}
El condensador es un circuito abierto en continua ($\abs{Z_C}\to\infty$ si $\omega\to0$) y un
cortocircuito a alta frecuencia ($\abs{Z_C}\to0$ si $\omega\to\infty$).

% -------------------------------------------------------------------------
\subsection{Transformada de Laplace y función de transferencia}\label{sec:transformada}
% -------------------------------------------------------------------------
Los fasores solo describen el régimen sinusoidal permanente. La transformada de Laplace
generaliza la idea a señales arbitrarias (cuadradas, triangulares, escalones), que son las que
se usan en el laboratorio para comprobar la integración y la derivación.

\begin{definicion}[Transformada de Laplace]\label{def:laplace}
Para una señal $f(t)$ definida en $t\ge0$, su transformada de Laplace es
\begin{equation}\label{eq:laplace}
  F(s)=\Lt\{f\}(s)=\int_{0}^{\infty}f(t)\,\eul^{-st}\,\dif t,\qquad s=\sigma+\jj\omega\in\C,
\end{equation}
definida en el semiplano $\Real s>\sigma_0$ en que la integral converge.
\end{definicion}

\begin{proposicion}[Propiedades que usaremos]\label{prop:laplace}
Si $f\leftrightarrow F$ y $g\leftrightarrow G$, con $a,b$ constantes:
\begin{gather}
  af(t)+bg(t)\;\longleftrightarrow\;aF(s)+bG(s),\label{eq:L-lineal}\\
  \frac{\dif f}{\dif t}\;\longleftrightarrow\;sF(s)-f(0),\qquad
  \int_{0}^{t}f(\tau)\,\dif\tau\;\longleftrightarrow\;\frac{F(s)}{s}.\label{eq:L-derivada}
\end{gather}
\end{proposicion}
\begin{proof}
La linealidad es la de la integral. Para la derivada se integra por partes:
$\int_0^\infty f'\eul^{-st}\dif t=\bigl[f\eul^{-st}\bigr]_0^\infty+s\int_0^\infty f\eul^{-st}\dif t=sF(s)-f(0)$.
Si $g(t)=\int_0^t f$, entonces $g'=f$ y $g(0)=0$, de modo que $F=sG$.
\end{proof}

Con condiciones iniciales nulas, \cref{eq:L-derivada} dice que \emph{derivar equivale a
multiplicar por $s$ e integrar, a dividir por $s$}. Aplicado al condensador,
$i=C\,\dif v/\dif t\Rightarrow I(s)=sC\,V(s)$, y se recupera \cref{eq:impedancias} con $\jj\omega\to s$:
\begin{equation}\label{eq:Zs}
  Z_R(s)=R,\qquad Z_C(s)=\frac{1}{sC},\qquad Z_L(s)=sL .
\end{equation}
Así, un circuito con condensadores se resuelve, en el dominio $s$, exactamente como uno resistivo.

\begin{definicion}[Función de transferencia]\label{def:transferencia}
Para un circuito lineal con condiciones iniciales nulas, la función de transferencia es el cociente
\begin{equation}\label{eq:transferencia}
  H(s)=\frac{V_{\mathrm{out}}(s)}{V_{\mathrm{in}}(s)} .
\end{equation}
Su restricción al eje imaginario, $H(\jj\omega)$, es la \emph{respuesta en frecuencia}.
\end{definicion}

\begin{proposicion}[Respuesta a una sinusoide]\label{prop:sinusoide}
Si $H(s)$ es estable (polos en $\Real s<0$) y $\vin(t)=V_p\cos\omega t$, la salida en régimen
permanente es
\begin{equation}\label{eq:respuesta-sinusoide}
  \vout(t)=\abs{H(\jj\omega)}\,V_p\cos\bigl(\omega t+\arg H(\jj\omega)\bigr).
\end{equation}
\end{proposicion}
\begin{proof}
La exponencial $\eul^{\jj\omega t}$ es autofunción de todo sistema lineal e invariante: si
$\vin=\eul^{\jj\omega t}$, la salida forzada es $H(\jj\omega)\eul^{\jj\omega t}$ (los términos asociados a los
polos de $H$ decaen como $\eul^{\Real(p)t}\to0$). Tomando la parte real, por linealidad, se obtiene el
resultado.
\end{proof}
La ganancia medida con el osciloscopio para una sinusoide de frecuencia $f$ es, por tanto,
$\abs{A_V}=\abs{H(\jj 2\pi f)}$, y el desfase entre los canales es $\arg H(\jj2\pi f)$.

% -------------------------------------------------------------------------
\subsection{Diagrama de Bode y red de constante temporal única}\label{sec:bode}
% -------------------------------------------------------------------------
\begin{definicion}[Diagrama de Bode]\label{def:bode}
Es la representación de $\abs{H(\jj\omega)}$ en decibelios, $20\log_{10}\abs{H}$, y de la fase
$\arg H(\jj\omega)$, ambas frente a la frecuencia en escala logarítmica. Una \emph{década} es un
factor $10$ en frecuencia y una \emph{octava}, un factor $2$.
\end{definicion}

\begin{proposicion}[Red de constante temporal única paso bajo]\label{prop:STC}
Sea
\begin{equation}\label{eq:STC}
  H(s)=\frac{K}{1+s/\omega_0},\qquad \omega_0=\frac1\tau .
\end{equation}
Entonces
\begin{equation}\label{eq:STC-modulo}
  \abs{H(\jj\omega)}=\frac{\abs{K}}{\sqrt{1+(\omega/\omega_0)^{2}}},\qquad
  \arg H(\jj\omega)=\arg K-\arctg\frac{\omega}{\omega_0},
\end{equation}
y el diagrama de Bode tiene dos asíntotas que se cortan en $\omega_0$ (\cref{fig:bode-stc}):
\begin{equation}\label{eq:STC-asintotas}
  \omega\ll\omega_0:\ \abs{H}\simeq\abs{K};\qquad
  \omega\gg\omega_0:\ \abs{H}\simeq\abs{K}\,\frac{\omega_0}{\omega},
  \ \text{pendiente }-20\dB/\text{década}\ (-6\dB/\text{octava}).
\end{equation}
En $\omega=\omega_0$ el módulo es $\abs{K}/\sqrt2$, es decir, $3\dB$ por debajo de la meseta, y la fase
se ha retrasado $45^{\circ}$: $\omega_0$ es la \emph{frecuencia de corte} a $-3\dB$.
\end{proposicion}
\begin{proof}
Con $s=\jj\omega$, $\abs{1+\jj x}=\sqrt{1+x^{2}}$ y $\arg(1+\jj x)=\arctg x$, con $x=\omega/\omega_0$. Si $x\ll1$,
$\sqrt{1+x^{2}}\simeq1$; si $x\gg1$, $\sqrt{1+x^{2}}\simeq x$, y $20\log_{10}(\abs{K}/x)$ disminuye
$20\dB$ cada vez que $x$ se multiplica por $10$ ($20\log_{10}2\simeq6\dB$ por octava). En $x=1$,
$20\log_{10}(1/\sqrt2)=-3{,}01\dB$ y $\arctg1=45^{\circ}$.
\end{proof}

Análogamente, $H(s)=K\,(s/\omega_0)/(1+s/\omega_0)$ es un \emph{paso alto}: crece a
$+20\dB$/década por debajo de $\omega_0$ y es plano por encima. Los dos casos límite
$H=K\omega_0/s$ (integrador, recta de $-20\dB$/década en todo el eje) y $H=Ks/\omega_0$
(derivador, $+20\dB$/década) son las asíntotas de alta y de baja frecuencia, respectivamente,
de estas dos redes.

% Figura: diagrama de Bode normalizado de una red STC paso bajo (K=1).
\begin{figure}[ht]
\centering
\begin{tikzpicture}
  \begin{semilogxaxis}[tema, name=mod, width=0.47\textwidth, height=5.2cm,
      xmin=0.01, xmax=100, ymin=-42, ymax=6, ytick={-40,-30,-20,-10,-3,0},
      xlabel={$\omega/\omega_0$}, ylabel={$\abs{H}/\abs{K}$ (dB)},
      log basis x=10, title style={font=\small}, title={(a) Módulo}]
    \addplot[asintota, domain=0.01:1] {0};
    \addplot[asintota, domain=1:100, samples=2] {-20*log10(x)};
    \addplot[azulacero, line width=1.1pt, domain=0.01:100, samples=200] {-10*log10(1+x^2)};
    \addplot[only marks, mark=*, mark size=1.6pt, rojoteo] coordinates {(1,-3.01)};
    \node[font=\scriptsize, anchor=south west, text=rojoteo] at (axis cs:1.15,-3) {$-3\dB$};
    \node[font=\scriptsize, anchor=south west, text=grispizarra, rotate=-33] at (axis cs:4,-11) {$-20\dB$/década};
  \end{semilogxaxis}
  \begin{semilogxaxis}[tema, at={(mod.east)}, anchor=west, xshift=1.9cm,
      width=0.47\textwidth, height=5.2cm,
      xmin=0.01, xmax=100, ymin=-95, ymax=5, ytick={-90,-45,0},
      xlabel={$\omega/\omega_0$}, ylabel={$\arg H-\arg K$ (grados)},
      log basis x=10, title style={font=\small}, title={(b) Fase}]
    \addplot[asintota, domain=0.01:0.1] {0};
    \addplot[asintota, domain=0.1:10, samples=2] {-45-45*log10(x)};
    \addplot[asintota, domain=10:100] {-90};
    \addplot[azulacero, line width=1.1pt, domain=0.01:100, samples=200] {-atan(x)};
    \addplot[only marks, mark=*, mark size=1.6pt, rojoteo] coordinates {(1,-45)};
    \node[font=\scriptsize, anchor=south west, text=rojoteo] at (axis cs:1.15,-44) {$-45^{\circ}$};
  \end{semilogxaxis}
\end{tikzpicture}
\caption{Diagrama de Bode de la red de constante temporal única paso bajo
$H(s)=K/(1+s/\omega_0)$, normalizado. En trazo discontinuo, las asíntotas de
\cref{eq:STC-asintotas} (en la fase, la aproximación usual de $-45^{\circ}$/década entre
$\omega_0/10$ y $10\,\omega_0$). En la frecuencia de corte el módulo cae $3\dB$ y la fase se retrasa
$45^{\circ}$.}
\label{fig:bode-stc}
\end{figure}


% =========================================================================
\section{El amplificador operacional}\label{sec:ao}
% =========================================================================
% -------------------------------------------------------------------------
\subsection{Definición, terminales y circuito equivalente}\label{sec:ao-def}
% -------------------------------------------------------------------------
\begin{definicion}[Amplificador operacional]\label{def:ao}
Un amplificador operacional es un amplificador de tensión con dos entradas, la \emph{inversora}
($-$) y la \emph{no inversora} ($+$), y una salida referida a tierra, que en su zona lineal
amplifica la diferencia de las tensiones de entrada:
\begin{equation}\label{eq:ao}
  \vout=A\,(\vp-\vm),
\end{equation}
donde $A$ es la \emph{ganancia en lazo abierto}. Se alimenta con dos fuentes continuas simétricas,
$+\Vcc$ y $-\Vee$, respecto de la tierra común (en el laboratorio, $\pm10\,\mathrm{V}$).
\end{definicion}

El circuito equivalente de la \figpanel{fig:ao}{b} es el modelo de cuadripolo de la
\cref{sec:cuadripolo} con entrada diferencial: una resistencia de entrada $\Rin$ entre las dos
entradas, y un equivalente de Thévenin $A(\vp-\vm)$ en serie con $\Rout$ a la salida~\cite{TEORIA-EA,SEDRA}.
En el circuito integrado $\mu$A741~\cite{UA741} las patillas son: $2$ (entrada $-$), $3$ (entrada $+$),
$6$ (salida), $7$ ($+\Vcc$), $4$ ($-\Vee$); la $1$ y la $5$ sirven para el ajuste de \emph{offset} y la $8$ no
se conecta. Por simplicidad, los esquemas omiten la alimentación, pero el AO \emph{siempre} está
alimentado: la corriente que entrega a la carga sale de las fuentes, no de la señal.

% Figura: (a) símbolo del AO con alimentación, (b) circuito equivalente, (c) característica en lazo abierto.
\begin{figure}[ht]
\centering
\begin{minipage}[b]{0.26\textwidth}\centering
\begin{circuitikz}[circuito, scale=0.9, transform shape]
  \draw (0,0) node[op amp, ao] (A) {};
  \draw (A.-) to[short, -o] ++(-0.6,0) node[left] {$\vm$};
  \draw (A.+) to[short, -o] ++(-0.6,0) node[left] {$\vp$};
  \draw (A.out) to[short, -o] ++(0.5,0) node[right] {$\vout$};
  \draw (A.up) -- ++(0,0.45) node[above, font=\scriptsize] {$+\Vcc$};
  \draw (A.down) -- ++(0,-0.45) node[below, font=\scriptsize] {$-\Vee$};
  \node[nota] at (-0.95,0.75) {2};
  \node[nota] at (-0.95,-0.75) {3};
  \node[nota] at (0.25,0.75) {7};
  \node[nota] at (0.25,-0.75) {4};
  \node[nota] at (1.15,0.2) {6};
\end{circuitikz}\\[2pt]
{\small (a)}
\end{minipage}\hfill
\begin{minipage}[b]{0.38\textwidth}\centering
\begin{circuitikz}[circuito, scale=0.8, transform shape]
  \draw[azulclaro, fill=opfondo, line width=0.5pt] (-0.4,2.8) -- (-0.4,-0.8) -- (4.9,1.0) -- cycle;
  \draw (-1.1,2.0) node[left] {$\vm$} to[short, o-] (0.6,2.0) to[R, l=$\Rin$] (0.6,0) -- (-1.1,0) to[short, -o] (-1.1,0) node[left] {$\vp$};
  \draw (2.1,-1.1) node[ground] {} -- (2.1,0.1) to[cV, l_=\scriptsize$A(\vp-\vm)$] (2.1,1.9) to[R, l=$\Rout$] (4.4,1.9) to[short, -o] (5.2,1.9) node[right] {$\vout$};
\end{circuitikz}\\[2pt]
{\small (b)}
\end{minipage}\hfill
\begin{minipage}[b]{0.30\textwidth}\centering
\begin{tikzpicture}
  \begin{axis}[temacentro, width=1.15\linewidth, height=4.2cm,
      xmin=-1.6, xmax=1.6, ymin=-1.45, ymax=1.45,
      xtick=\empty, ytick=\empty,
      xlabel={$\vp-\vm$},
      xlabel style={at={(axis description cs:1,0.5)}, anchor=south east, yshift=2pt, font=\scriptsize}]
    \addplot[senal salida, samples=3, domain=-1.6:-0.35] {-1};
    \addplot[senal salida, samples=2, domain=-0.35:0.35] {x/0.35};
    \addplot[senal salida, samples=3, domain=0.35:1.6] {1};
    \addplot[punteado] coordinates {(0.35,0) (0.35,1)};
    \addplot[punteado] coordinates {(-0.35,0) (-0.35,-1)};
    \node[font=\scriptsize, anchor=south] at (axis cs:1.0,1.0) {$+\Vsat$};
    \node[font=\scriptsize, anchor=north] at (axis cs:-1.0,-1.0) {$-\Vsat$};
    \node[font=\scriptsize, anchor=north] at (axis cs:0.42,0.0) {$\frac{\Vsat}{A}$};
    \node[font=\scriptsize, anchor=east, text=rojoteo] at (axis cs:-0.08,0.6) {pendiente $A$};
    \node[font=\scriptsize, anchor=west] at (axis description cs:0.52,0.97) {$\vout$};
  \end{axis}
\end{tikzpicture}\\[2pt]
{\small (c)}
\end{minipage}
\caption{El amplificador operacional. (a) Símbolo, con las patillas del $\mu$A741 y la
alimentación (que en los esquemas siguientes se omite). (b) Circuito equivalente: resistencia
de entrada diferencial $\Rin$ y equivalente de Thévenin de salida~\cite{TEORIA-EA,SEDRA}. (c)
Característica de transferencia en lazo abierto (escala horizontal muy ampliada: la zona lineal
mide $2\Vsat/A\sim\SI{0.1}{\milli\volt}$).}
\label{fig:ao}
\end{figure}


% -------------------------------------------------------------------------
\subsection{Saturación}\label{sec:saturacion}
% -------------------------------------------------------------------------
La salida no puede exceder las tensiones de alimentación. En un AO como el 741 queda, además,
un margen de en torno a $1$--$2\,\mathrm{V}$ con cada fuente~\cite{UA741}:
\begin{equation}\label{eq:saturacion}
  -\Vsat\le\vout\le+\Vsat,\qquad \Vsat\lesssim\Vcc-1\,\mathrm{V}
  \quad(\Vsat\approx9\,\mathrm{V}\ \text{para}\ \Vcc=10\,\mathrm{V}).
\end{equation}
La característica de transferencia en lazo abierto, \cref{eq:ao} con \cref{eq:saturacion}, es
la de la \figpanel{fig:ao}{c}: una recta de pendiente $A$ limitada por $\pm\Vsat$. Como $A\sim10^{5}$,
la zona lineal es solo $\abs{\vp-\vm}<\Vsat/A\sim\SI{90}{\micro\volt}$: sin realimentación, el AO
se comporta como un comparador. Valores concretos del 741 se recogen en la \cref{tab:741}.

% -------------------------------------------------------------------------
\subsection{Modelo ideal y realimentación negativa}\label{sec:ideal}
% -------------------------------------------------------------------------
\begin{hipotesis}[Amplificador operacional ideal]\label{hip:ideal}
\begin{equation}\label{eq:ideal}
  \Rin\to\infty,\qquad \Rout\to0,\qquad A\to\infty\quad(\text{e independiente de la frecuencia}).
\end{equation}
\end{hipotesis}

\begin{proposicion}[Reglas del AO ideal]\label{prop:reglas}
Si el AO ideal trabaja en su zona lineal, gracias a una realimentación negativa (una red que
une la salida con la entrada inversora), se cumplen:
\begin{align}
  &\text{(R1) corrientes de entrada nulas:} && i_{+}=i_{-}=0, \label{eq:R1}\\
  &\text{(R2) cortocircuito virtual:} && \vp=\vm . \label{eq:R2}
\end{align}
Además, por $\Rout=0$, la tensión de salida no depende de la carga conectada.
\end{proposicion}
\begin{proof}
(R1) Entre las entradas solo hay $\Rin$; con $\Rin\to\infty$ no circula corriente por ella.
(R2) En zona lineal, $\vout$ es finita ($\abs{\vout}\le\Vsat$), luego por \cref{eq:ao}
$\vp-\vm=\vout/A\to0$ si $A\to\infty$. La realimentación negativa es la que hace que esta situación
sea estable: si $\vout$ crece, $\vm$ crece con ella, $\vp-\vm$ disminuye y $\vout$ se corrige.
Con la salida en vacío, $\Rout=0$ hace que el divisor de salida de \cref{eq:carga} valga $1$.
\end{proof}

\begin{observacion}[Tierra virtual]\label{obs:tierra-virtual}
Si $\vp=0$ (entrada no inversora a tierra), (R2) da $\vm=0$: la entrada inversora está a potencial
de tierra sin estar conectada a ella. Es una \emph{tierra virtual}: tiene tensión nula pero no absorbe
corriente, por (R1). Las dos reglas son incompatibles con la saturación: si $\abs{\vout}=\Vsat$, el lazo
ya no actúa y deja de cumplirse \cref{eq:R2}.
\end{observacion}

\begin{proposicion}[Método de análisis]\label{prop:metodo}
Todo circuito con AO ideal y realimentación negativa se resuelve con cuatro pasos:
\begin{enumerate}
  \item obtener $\vp$ a partir del circuito conectado a la entrada no inversora (por (R1), sin carga);
  \item imponer $\vm=\vp$ (R2);
  \item aplicar la ley de Kirchhoff de corrientes en el nudo inversor con $i_{-}=0$ (R1), usando
        impedancias si hay elementos reactivos;
  \item despejar $\vout$ y comprobar que $\abs{\vout}<\Vsat$ (si no, el resultado no es válido).
\end{enumerate}
\end{proposicion}

% -------------------------------------------------------------------------
\subsection{Límites del modelo ideal}\label{sec:real}
% -------------------------------------------------------------------------
El modelo ideal es excelente a baja frecuencia y con resistencias de valor moderado
($1$--$100\,\mathrm{k}\Omega$). Sus límites explican lo que se observa al forzar el circuito
(amplitudes grandes o frecuencias de $100\,\mathrm{kHz}$).

\begin{proposicion}[Ganancia en lazo abierto real]\label{prop:Areal}
La ganancia del AO tiene un polo dominante a una frecuencia $f_b$ muy baja~\cite{TEORIA-EA}:
\begin{equation}\label{eq:Areal}
  A(s)=\frac{A_0}{1+s/\omega_b},\qquad
  \abs{A(\jj\omega)}\simeq\frac{A_0\,\omega_b}{\omega}=\frac{\omega_t}{\omega}\quad(\omega\gg\omega_b),\qquad
  f_t=A_0f_b ,
\end{equation}
donde $f_t$, la frecuencia de ganancia unidad, es el \emph{producto ganancia--ancho de banda} (GBW).
\end{proposicion}
\begin{proof}
Es la red STC de la \cref{prop:STC} con $K=A_0$; por encima de $\omega_b$, $\abs{A}=A_0\omega_b/\omega$
se anula ($0\dB$) en $\omega_t=A_0\omega_b$.
\end{proof}

\begin{proposicion}[Ancho de banda en lazo cerrado]\label{prop:BW}
Con ganancia finita $A$, los amplificadores inversor y no inversor de las \cref{sec:pl-inversor,sec:pl-no-inversor} tienen
\begin{equation}\label{eq:lazo-cerrado}
  A_V^{\mathrm{inv}}=\frac{-R_2/R_1}{1+\dfrac{1+R_2/R_1}{A}},\qquad
  A_V^{\mathrm{no\,inv}}=\frac{1+R_2/R_1}{1+\dfrac{1+R_2/R_1}{A}},
\end{equation}
que tienden a los valores ideales si $A\gg1+R_2/R_1$. Con \cref{eq:Areal}, ambos son redes STC paso bajo de frecuencia de corte
\begin{equation}\label{eq:f3dB}
  f_{3\mathrm{dB}}=\frac{f_t}{1+R_2/R_1} .
\end{equation}
\end{proposicion}
\begin{proof}
Con $i_-=0$, el nudo inversor es un divisor entre la entrada y la salida. Para el inversor,
por superposición, $\vm=\vin\,R_2/(R_1+R_2)+\vout\,R_1/(R_1+R_2)$, y $\vout=-A\vm$. Llamando
$\beta=R_1/(R_1+R_2)$, $\vout(1+A\beta)=-A(1-\beta)\vin$, y dividiendo por $A\beta$ se obtiene la
primera expresión, pues $(1-\beta)/\beta=R_2/R_1$ y $1/\beta=1+R_2/R_1$. Para el no inversor,
$\vm=\beta\vout$ y $\vout=A(\vin-\beta\vout)$, de donde $\vout/\vin=(1/\beta)/(1+1/(A\beta))$. Si
$A=\omega_t/s$, el denominador común es $1+s/(\beta\omega_t)$: un polo en $\beta\omega_t$.
\end{proof}

Es decir, ganancia y ancho de banda se intercambian: el producto $\abs{A_V}\cdot f_{3\mathrm{dB}}\approx f_t$ es
constante (\cref{fig:gbw}). Las otras limitaciones relevantes son:
\begin{itemize}
  \item \emph{Velocidad de subida} (\emph{slew rate}): $\abs{\dif\vout/\dif t}\le\mathrm{SR}$. Una sinusoide
        $\vout=V_p\sen\omega t$ tiene pendiente máxima $\omega V_p$; para no deformarse en un
        triángulo debe cumplir
        \begin{equation}\label{eq:slew}
          2\pi f\,V_p\le\mathrm{SR}.
        \end{equation}
  \item \emph{Corriente de salida limitada} (cortocircuito, $\sim\SI{25}{\milli\ampere}$): la salida no
        puede alimentar cargas de pocos ohmios.
  \item \emph{Corrientes de polarización} $I_B$ y \emph{tensión de offset} $V_{IO}$: no nulas; su efecto en
        la salida crece con el valor de las resistencias (o, en un integrador, con el tiempo).
\end{itemize}

% Figura: Bode del AO en lazo abierto (A0=2e5, fb=5 Hz, ft=1 MHz) y de las ganancias en lazo cerrado.
\begin{figure}[ht]
\centering
\begin{tikzpicture}
  \begin{semilogxaxis}[tema, width=0.82\textwidth, height=5.8cm,
      xmin=1, xmax=1e7, ymin=-5, ymax=115, ytick={0,20,40,60,80,100},
      xlabel={$f$ (Hz)}, ylabel={Módulo (dB)}, log basis x=10,
      legend pos=north east]
    \addplot[azulacero, line width=1.1pt, domain=1:1e7, samples=300]
      {20*log10(2e5/sqrt(1+(x/5)^2))};
    \addlegendentry{$\abs{A(\jj\omega)}$, lazo abierto}
    \addplot[rojoteo, line width=1.1pt, domain=1:1e7, samples=300]
      {20*log10(10/sqrt(1+(x/90909)^2))};
    \addlegendentry{$\abs{A_V}$, inversor ($R_2/R_1=10$)}
    \addplot[asintota, domain=1:1e7, samples=2] {20};
    \addplot[punteado] coordinates {(90909,-5) (90909,17)};
    \addplot[punteado] coordinates {(5,-5) (5,103)};
    \addplot[only marks, mark=*, mark size=1.5pt, rojoteo] coordinates {(90909,17)};
    \addplot[only marks, mark=*, mark size=1.5pt, azulacero] coordinates {(5,103) (1e6,0)};
    \node[font=\scriptsize, anchor=south west] at (axis cs:6,103) {$f_b\approx5\,\mathrm{Hz}$};
    \node[font=\scriptsize, anchor=south west] at (axis cs:1.1e6,1) {$f_t\approx1\,\mathrm{MHz}$};
    \node[font=\scriptsize, anchor=west, text=rojoteo] (f3) at (axis cs:150,8) {$f_{3\mathrm{dB}}=f_t/(1+R_2/R_1)\approx91\,\mathrm{kHz}$};
    \draw[->, rojoteo, line width=0.4pt] (f3.east) -- (axis cs:80000,15.5);
    \node[font=\scriptsize, anchor=south, text=grispizarra, rotate=-27] at (axis cs:3000,58) {$-20\dB$/década};
  \end{semilogxaxis}
\end{tikzpicture}
\caption{Diagrama de Bode (módulo) de la ganancia en lazo abierto de un AO de tipo 741, según
\cref{eq:Areal} con $A_0=2\cdot10^{5}$ y $f_t=1\,\mathrm{MHz}$, y de la ganancia en lazo cerrado del
amplificador inversor de la \cref{sec:pl-inversor}, según \cref{eq:lazo-cerrado,eq:f3dB}. La ganancia
en lazo cerrado sigue la recta de $20\dB$ hasta cortar la de lazo abierto: a partir de ahí manda
el AO.}
\label{fig:gbw}
\end{figure}


\begin{table}[ht]
\centering\small
\renewcommand{\arraystretch}{1.2}
\caption{Parámetros típicos del $\mu$A741 a $25\,^{\circ}\mathrm{C}$~\cite{UA741} frente al modelo ideal.}
\label{tab:741}
\begin{tabular}{@{}llcc@{}}
\toprule
Parámetro & Símbolo & $\mu$A741 (típico) & Ideal\\
\midrule
Ganancia en lazo abierto (continua) & $A_0$ & $2\cdot10^{5}$ ($106\dB$) & $\infty$\\
Producto ganancia--ancho de banda & $f_t$ & $\approx\SI{1}{\mega\hertz}$ & $\infty$\\
Resistencia de entrada & $\Rin$ & $\SI{2}{\mega\ohm}$ & $\infty$\\
Resistencia de salida & $\Rout$ & $\SI{75}{\ohm}$ & $0$\\
Velocidad de subida & SR & $\SI{0.5}{\volt\per\micro\second}$ & $\infty$\\
Corriente de polarización & $I_B$ & $\SI{80}{\nano\ampere}$ & $0$\\
Tensión de \emph{offset} de entrada & $V_{IO}$ & $\SI{1}{\milli\volt}$ & $0$\\
Corriente de salida en cortocircuito & $I_{OS}$ & $\pm\SI{25}{\milli\ampere}$ & $\infty$\\
Excursión de salida ($R_L\ge\SI{10}{\kilo\ohm}$) & $\pm\Vsat$ & $\approx\pm(\Vcc-1\,\mathrm{V})$ & $\pm\Vcc$\\
\bottomrule
\end{tabular}
\end{table}


% =========================================================================
\parte{II}{PreLab}
% =========================================================================
% =========================================================================
\section{Planteamiento del PreLab}\label{sec:pl-planteamiento}
% =========================================================================
El PreLab es la fase de diseño de la sesión~\cite{GUION-S1}: con el modelo ideal del AO se obtienen
las ganancias de los tres circuitos en función de sus componentes y se eligen los valores que dan
las ganancias pedidas. Los recuadros amarillos del guion son tres; cada uno se reproduce como
\emph{PreLab~$n$} y va seguido de su resolución. Todos los cálculos comparten las hipótesis siguientes.

\begin{hipotesis}[Hipótesis comunes]\label{hip:comunes}
Salvo que se indique otra cosa, en todos los cálculos se supone que:
\begin{enumerate}
  \item El AO ($\mu$A741) es ideal (\cref{hip:ideal}) y trabaja en zona lineal gracias a la
        realimentación negativa, de modo que valen las reglas (R1) y (R2) de la \cref{prop:reglas}.
  \item Se alimenta con $\pm\Vcc=\pm10\,\mathrm{V}$, de modo que $\Vsat\approx9\,\mathrm{V}$ (\cref{eq:saturacion}).
  \item La señal de entrada es sinusoidal de $f=1\,\mathrm{kHz}$ y el circuito está en régimen permanente
        (en el circuito~3, frecuencia variable).
  \item La carga de salida es $R_L=10\,\mathrm{k}\Omega$ y el generador tiene resistencia interna $R_s=50\,\Omega$.
\end{enumerate}
\end{hipotesis}

\begin{observacion}[Cómo trasladar cada PreLab al manuscrito]\label{obs:manuscrito}
Para cada recuadro: (i) copiar el enunciado; (ii) dibujar el esquema del circuito con los nudos,
las tensiones y los sentidos de las corrientes que se usan en el cálculo; (iii) escribir cada paso
con la regla que lo justifica (R1, R2, Kirchhoff, Ohm, impedancia), con las ecuaciones numeradas a la
derecha; (iv) enmarcar el resultado general y el valor numérico del diseño; (v) añadir las gráficas,
con los ejes rotulados y las unidades. Los esquemas y gráficas de este documento son los que deben
reproducirse a mano.
\end{observacion}

% =========================================================================
\section{Circuito 1: amplificador inversor}\label{sec:pl-inversor}
% =========================================================================
\begin{prelab}[Amplificador inversor]\label{pl:inversor}
Usando el modelo de amplificador operacional ideal visto en teoría, encuéntrese la ganancia en
tensión $A_V=\Vout/\Vin$ en función de las resistencias del circuito (\cref{fig:inversor}). Si
imponemos $R_1=1\,\mathrm{k}\Omega$, determínese el valor de $R_2$ para conseguir $A_V=-10$ a $1\,\mathrm{kHz}$.
\end{prelab}

% Figura: amplificador inversor (circuito 1) con tensiones y corrientes del análisis.
\begin{figure}[ht]
\centering
\begin{circuitikz}[circuito]
  \draw (0,0) node[op amp, ao] (A) {};
  % nudo inversor y red de realimentación
  \draw (A.-) -- ++(-0.6,0) coordinate (N);
  \draw (N) to[R, l_=$R_1$, i<^=$i_1$] ++(-2.4,0) coordinate (E);
  \draw (N) |- ++(0,1.5) coordinate (T);
  \draw (A.out) -- ++(0.7,0) coordinate (O);
  \draw (T) to[R, l=$R_2$, i>^=$i_2$] (T -| O) -- (O);
  % generador
  \draw (E) to[short, -o] ++(-0.3,0) coordinate (Ein) node[above, tens] {$\vin$};
  \draw (Ein) -- ++(-0.9,0) coordinate (G) to[sV, l_=$v_g$] ++(0,-2.0) node[ground] {};
  % entrada no inversora a tierra
  \draw (A.+) -- ++(-0.6,0) coordinate (P) -- ++(0,-0.45) node[ground] {};
  % carga y salida
  \draw (O) to[short, *-o] ++(1.0,0) coordinate (S) node[right, tens] {$\vout$};
  \draw (O) ++(0.5,0) to[R, l=$R_L$] ++(0,-2.0) node[ground] {};
  \fill (N) circle (1.6pt);
  % anotaciones del análisis
  \node[tens, anchor=east, font=\scriptsize] at ($(N)+(-0.08,0.75)$) {$\vm=0$};
  \draw[azulprof, line width=0.3pt] ($(N)+(-0.08,0.72)$) -- ($(N)+(-0.02,0.06)$);
  \node[tens, anchor=east, font=\scriptsize] at ($(P)+(-0.1,-0.25)$) {$\vp=0$};
\end{circuitikz}
\caption{Amplificador inversor (figura~2 del guion~\cite{GUION-S1}), con $R_1=1\,\mathrm{k}\Omega$,
$R_L=10\,\mathrm{k}\Omega$ y la alimentación omitida. Se indican las magnitudes del análisis: el nudo
inversor es una tierra virtual ($\vm=\vp=0$) y no absorbe corriente, de modo que $i_1=i_2$.}
\label{fig:inversor}
\end{figure}


\begin{resolucion}
Con los sentidos de corriente de la \cref{fig:inversor}:

\textbf{1. Tierra virtual.} La entrada $+$ está a tierra; por (R2),
\begin{equation}\label{eq:inv-vm}
  \vm=\vp=0 .
\end{equation}

\textbf{2. Ley de Ohm en $R_1$.}
\begin{equation}\label{eq:inv-i1}
  i_1=\frac{\vin-\vm}{R_1}=\frac{\vin}{R_1} .
\end{equation}

\textbf{3. Kirchhoff en el nudo inversor}, con $i_-=0$ por (R1):
\begin{equation}\label{eq:inv-kcl}
  i_2=i_1-i_-=\frac{\vin}{R_1} .
\end{equation}

\textbf{4. Ley de Ohm en $R_2$} ($i_2$ va del nudo inversor a la salida):
\begin{equation}\label{eq:inv-vout}
  \vout=\vm-i_2R_2=-\frac{R_2}{R_1}\,\vin .
\end{equation}

\textbf{Resultado.}
\begin{equation}\label{eq:inv-Av}
  \resultado{A_V=\frac{\vout}{\vin}=-\frac{R_2}{R_1}}
\end{equation}

\textbf{Diseño.} Con $A_V=-10$ y $R_1=1\,\mathrm{k}\Omega$ en \cref{eq:inv-Av}:
\begin{equation}\label{eq:inv-R2}
  R_2=\abs{A_V}\,R_1\quad\Longrightarrow\quad\resultado{R_2=10\,\mathrm{k}\Omega}
\end{equation}

\textbf{Comentario.} El signo indica contrafase ($180^{\circ}$). En el modelo ideal $A_V$ no depende de la
frecuencia (vale igual a $1\,\mathrm{kHz}$), ni de $A$ ni de $R_L$: solo del cociente de resistencias
(\cref{fig:inv-graficas}).
\end{resolucion}

% Figura: previsiones del amplificador inversor (A_V=-10, Vsat=9 V, f=1 kHz).
\begin{figure}[ht]
\centering
\begin{tikzpicture}
  % (a) zona lineal
  \begin{axis}[tema, name=a, width=0.38\textwidth, height=4.6cm,
      xmin=0, xmax=2, ymin=-10.5, ymax=10.5, ytick={-9,-5,0,5,9},
      xlabel={$t$ (ms)}, ylabel={Tensión (V)}, title style={font=\small},
      title={(a) $V_p=0{,}5\,\mathrm{V}$}]
    \addplot[asintota, domain=0:2, samples=2] {9};
    \addplot[asintota, domain=0:2, samples=2] {-9};
    \addplot[senal entrada, domain=0:2, samples=150] {0.5*sin(360*x)};
    \addplot[senal salida, domain=0:2, samples=150] {-5*sin(360*x)};
  \end{axis}
  % (b) característica de transferencia
  \begin{axis}[temacentro, name=b, at={(a.east)}, anchor=west, xshift=0.9cm,
      width=0.33\textwidth, height=4.6cm,
      xmin=-1.6, xmax=1.6, ymin=-11, ymax=11, xtick={-0.9,0.9}, ytick={-9,9},
      xticklabels={$-0{,}9$,$0{,}9$}, yticklabels={$-9$,$9$},
      xlabel={$\vin$ (V)}, 
      xlabel style={at={(axis description cs:1,0.5)}, anchor=south east, yshift=2pt, font=\scriptsize}, clip=false,
      title style={font=\small, yshift=10pt}, title={(b) $\vout(\vin)$}]
    \addplot[senal salida, domain=-1.6:1.6, samples=200] {max(-9,min(9,-10*x))};
    \addplot[punteado] coordinates {(0.9,0) (0.9,-9)};
    \addplot[punteado] coordinates {(-0.9,0) (-0.9,9)};
    \node[font=\scriptsize, anchor=west, text=rojoteo] at (axis cs:0.1,6.6) {pendiente $-10$};
    \node[font=\scriptsize, anchor=west] at (axis description cs:0.52,0.97) {$\vout$ (V)};
  \end{axis}
  % (c) recorte
  \begin{axis}[tema, at={(b.east)}, anchor=west, xshift=0.7cm, width=0.38\textwidth, height=4.6cm,
      xmin=0, xmax=2, ymin=-16, ymax=16, ytick={-15,-9,0,9,15}, yticklabels={},
      xlabel={$t$ (ms)}, title style={font=\small},
      title={(c) $V_p=1{,}5\,\mathrm{V}$}]
    \addplot[gray!60, densely dashed, domain=0:2, samples=150] {-15*sin(360*x)};
    \addplot[asintota, domain=0:2, samples=2] {9};
    \addplot[asintota, domain=0:2, samples=2] {-9};
    \addplot[senal entrada, domain=0:2, samples=150] {1.5*sin(360*x)};
    \addplot[senal salida, domain=0:2, samples=300] {max(-9,min(9,-15*sin(360*x)))};
    \node[font=\scriptsize, anchor=south east] at (axis cs:2,9) {$+\Vsat$};
  \end{axis}
\end{tikzpicture}
\caption{Previsiones para el amplificador inversor ($A_V=-10$, $\Vsat\approx9\,\mathrm{V}$,
$f=1\,\mathrm{kHz}$); en azul la entrada y en ámbar la salida. (a) En zona lineal la salida es la
entrada multiplicada por $-10$ (contrafase). (b) Característica de transferencia: recta de pendiente
$-10$ hasta $\abs{\vin}=\Vsat/10=0{,}9\,\mathrm{V}$ y saturación fuera. (c) Con $V_p=1{,}5\,\mathrm{V}$ la salida
lineal, en gris, superaría $\Vsat$ y la salida real queda recortada en $\pm\Vsat$ (\emph{clipping}).}
\label{fig:inv-graficas}
\end{figure}


% =========================================================================
\section{Circuito 2: amplificador no inversor}\label{sec:pl-no-inversor}
% =========================================================================
\begin{prelab}[Amplificador no inversor]\label{pl:no-inversor}
Usando un modelo de amplificador operacional ideal, encuéntrese la ganancia en tensión
$A_V=\Vout/\Vin$ en función de las resistencias del circuito de la \cref{fig:no-inversor}. Si
imponemos $R_1=1\,\mathrm{k}\Omega$, determínese el valor de $R_2$ para conseguir $A_V=+11$.
\end{prelab}

% Figura: amplificador no inversor (circuito 2) con tensiones y corrientes del análisis.
\begin{figure}[ht]
\centering
\begin{circuitikz}[circuito]
  \draw (0,0) node[op amp, ao, noinv input up] (A) {};
  % nudo inversor (abajo) y red de realimentación
  \draw (A.out) -- ++(0.7,0) coordinate (O);
  \draw (A.-) -- ++(-0.6,0) coordinate (N) -- ++(0,-1.0) coordinate (M);
  \draw (M) -- ++(-1.0,0) coordinate (G1) to[R, l_=$R_1$, i>^=$i_1$] ++(0,-1.7) node[ground] {};
  \draw (M) to[R, l_=$R_2$, i<_=$i_2$] (M -| O) -- (O);
  % generador en la entrada no inversora
  \draw (A.+) to[short, -o] ++(-2.0,0) coordinate (Ein) node[above, tens] {$\vin$};
  \draw (Ein) -- ++(-0.8,0) to[sV, l_=$v_g$] ++(0,-2.0) node[ground] {};
  % carga y salida
  \draw (O) to[short, *-o] ++(1.0,0) node[right, tens] {$\vout$};
  \draw (O) ++(0.5,0) to[R, l=$R_L$] ++(0,-2.0) node[ground] {};
  \fill (M) circle (1.6pt);
  \node[tens, anchor=east, font=\scriptsize] at ($(N)!0.5!(M)+(-0.08,0)$) {$\vm=\vin$};
\end{circuitikz}
\caption{Amplificador no inversor (figura~3 del guion~\cite{GUION-S1}), con $R_1=1\,\mathrm{k}\Omega$,
$R_L=10\,\mathrm{k}\Omega$ y la alimentación omitida. La entrada $+$ no absorbe corriente; el nudo
inversor sigue a la entrada ($\vm=\vin$) y $R_1$--$R_2$ forman un divisor de $\vout$.}
\label{fig:no-inversor}
\end{figure}


\begin{resolucion}
Con los sentidos de corriente de la \cref{fig:no-inversor}:

\textbf{1. Cortocircuito virtual.} La señal entra por $+$ y, por (R1), no hay caída en esa conexión;
por (R2),
\begin{equation}\label{eq:ni-vm}
  \vm=\vp=\vin .
\end{equation}

\textbf{2. Ley de Ohm en $R_1$} (del nudo inversor a tierra):
\begin{equation}\label{eq:ni-i1}
  i_1=\frac{\vm-0}{R_1}=\frac{\vin}{R_1} .
\end{equation}

\textbf{3. Kirchhoff en el nudo inversor}, con $i_-=0$ por (R1):
\begin{equation}\label{eq:ni-kcl}
  i_2=i_1+i_-=\frac{\vin}{R_1} .
\end{equation}

\textbf{4. Ley de Ohm en $R_2$} ($i_2$ va de la salida al nudo inversor):
\begin{equation}\label{eq:ni-vout}
  \vout=\vm+i_2R_2=\Bigl(1+\frac{R_2}{R_1}\Bigr)\vin .
\end{equation}

\textbf{Resultado.}
\begin{equation}\label{eq:ni-Av}
  \resultado{A_V=\frac{\vout}{\vin}=1+\frac{R_2}{R_1}}
\end{equation}

\textbf{Diseño.} Con $A_V=11$ y $R_1=1\,\mathrm{k}\Omega$ en \cref{eq:ni-Av}:
\begin{equation}\label{eq:ni-R2}
  R_2=(A_V-1)\,R_1\quad\Longrightarrow\quad\resultado{R_2=10\,\mathrm{k}\Omega}
\end{equation}

\textbf{Comentario.} $A_V>0$: salida en fase con la entrada, y $A_V\ge1$ siempre. Como en el inversor,
la ganancia solo depende del cociente de resistencias (\cref{fig:ni-graficas}).
\end{resolucion}

% Figura: previsiones del amplificador no inversor (A_V=+11, Vsat=9 V, f=1 kHz).
\begin{figure}[ht]
\centering
\begin{tikzpicture}
  \begin{axis}[tema, name=a, width=0.45\textwidth, height=4.6cm,
      xmin=0, xmax=2, ymin=-10.5, ymax=10.5, ytick={-9,-5.5,0,5.5,9},
      xlabel={$t$ (ms)}, ylabel={Tensión (V)}, title style={font=\small},
      title={(a) $V_p=0{,}5\,\mathrm{V}$}]
    \addplot[asintota, domain=0:2, samples=2] {9};
    \addplot[asintota, domain=0:2, samples=2] {-9};
    \addplot[senal entrada, domain=0:2, samples=150] {0.5*sin(360*x)};
    \addplot[senal salida, domain=0:2, samples=150] {5.5*sin(360*x)};
  \end{axis}
  \begin{axis}[temacentro, at={(a.east)}, anchor=west, xshift=1.2cm,
      width=0.42\textwidth, height=4.6cm,
      xmin=-1.6, xmax=1.6, ymin=-11, ymax=11, xtick={-0.82,0.82}, ytick={-9,9},
      xticklabels={$-0{,}82$,$0{,}82$}, yticklabels={$-9$,$9$},
      xlabel={$\vin$ (V)}, 
      xlabel style={at={(axis description cs:1,0.5)}, anchor=south east, yshift=2pt, font=\scriptsize}, clip=false,
      title style={font=\small, yshift=10pt}, title={(b) $\vout(\vin)$}]
    \addplot[senal salida, domain=-1.6:1.6, samples=200] {max(-9,min(9,11*x))};
    \addplot[punteado] coordinates {(0.818,0) (0.818,9)};
    \addplot[punteado] coordinates {(-0.818,0) (-0.818,-9)};
    \node[font=\scriptsize, anchor=west, text=rojoteo] at (axis cs:-1.55,5.8) {pendiente $+11$};
    \node[font=\scriptsize, anchor=west] at (axis description cs:0.52,0.97) {$\vout$ (V)};
  \end{axis}
\end{tikzpicture}
\caption{Previsiones para el amplificador no inversor ($A_V=+11$, $\Vsat\approx9\,\mathrm{V}$,
$f=1\,\mathrm{kHz}$). (a) Entrada (azul) y salida (ámbar) en fase. (b) Característica de
transferencia: pendiente $+11$ hasta $\abs{\vin}=\Vsat/11\approx0{,}82\,\mathrm{V}$ y saturación fuera.}
\label{fig:ni-graficas}
\end{figure}


% =========================================================================
\section{Circuito 3: amplificador inversor generalizado y filtro paso bajo}\label{sec:pl-filtro}
% =========================================================================
\begin{prelab}[Inversor generalizado y filtro]\label{pl:filtro}
Sea el circuito de la \figpanel{fig:generalizado}{a}, en el que las impedancias $Z_1$ y $Z_2$ pueden
ser no resistivas y, por lo tanto, tener dependencia en frecuencia: obténgase la ganancia en voltaje
del circuito. Esta ganancia será función de la frecuencia salvo que ambas impedancias sean puramente
resistivas. Considérese el caso de que $Z_1=R$ sea real y $Z_2$ una composición en paralelo de una
resistencia $R_F$ y un condensador $C$ (\figpanel{fig:generalizado}{b}). Demuéstrese que la dependencia
de la función de transferencia con la frecuencia $f$ es la de un filtro paso bajo: para frecuencias
suficientemente bajas la ganancia en tensión es constante, mientras que para frecuencias
suficientemente altas disminuye.
\end{prelab}

% Figura: (a) amplificador inversor generalizado; (b) caso Z1=R, Z2=RF||C (filtro paso bajo).
\begin{figure}[ht]
\centering
\begin{minipage}[b]{0.46\textwidth}\centering
\begin{circuitikz}[circuito, scale=0.88, transform shape]
  \draw (0,0) node[op amp, ao] (A) {};
  \draw (A.-) -- ++(-0.5,0) coordinate (N);
  \draw (N) to[generic, l_=$Z_1$, i<^=$I_1$] ++(-2.2,0) to[short, -o] ++(-0.3,0) node[left, tens] {$\Vin$};
  \draw (A.out) -- ++(0.6,0) coordinate (O);
  \draw (N) -- ++(0,1.4) coordinate (T) to[generic, l=$Z_2$, i>^=$I_2$] (T -| O) -- (O);
  \draw (O) to[short, *-o] ++(0.6,0) node[right, tens] {$\Vout$};
  \draw (A.+) -- ++(-0.5,0) -- ++(0,-0.4) node[ground] {};
  \fill (N) circle (1.6pt);
\end{circuitikz}\\[2pt]
{\small (a)}
\end{minipage}\hfill
\begin{minipage}[b]{0.50\textwidth}\centering
\begin{circuitikz}[circuito, scale=0.88, transform shape]
  \draw (0,0) node[op amp, ao] (A) {};
  \draw (A.-) -- ++(-0.5,0) coordinate (N);
  \draw (N) to[R, l_=$R$] ++(-2.2,0) to[short, -o] ++(-0.3,0) node[left, tens] {$\Vin$};
  \draw (A.out) -- ++(0.6,0) coordinate (O);
  \draw (N) -- ++(0,1.2) coordinate (T1) to[C, l_=$C$] (T1 -| O) -- (O);
  \draw (T1) -- ++(0,1.0) coordinate (T2) to[R, l=$R_F$] (T2 -| O) -- (T1 -| O);
  \fill (T1) circle (1.6pt); \fill (T1 -| O) circle (1.6pt);
  \draw (O) to[short, *-o] ++(0.6,0) node[right, tens] {$\Vout$};
  \draw (A.+) -- ++(-0.5,0) -- ++(0,-0.4) node[ground] {};
  \fill (N) circle (1.6pt);
\end{circuitikz}\\[2pt]
{\small (b)}
\end{minipage}
\caption{(a) Amplificador inversor generalizado (figura~4a del guion~\cite{GUION-S1}): $Z_1(s)$ y $Z_2(s)$
son impedancias arbitrarias. (b) Caso $Z_1=R$, $Z_2=R_F\parallel C$ (figura~5a del guion), que es un
filtro paso bajo activo y, para $R_F\to\infty$, un integrador.}
\label{fig:generalizado}
\end{figure}


% -------------------------------------------------------------------------
\subsection{Ganancia del inversor generalizado}\label{sec:pl-generalizado}
% -------------------------------------------------------------------------
\begin{resolucion}
Trabajamos en el dominio de Laplace con condiciones iniciales nulas: cada elemento cumple
$V(s)=Z(s)\,I(s)$, \cref{eq:Zs}, y Kirchhoff, (R1) y (R2) se mantienen por linealidad. El cálculo es
el del inversor con $R_k\to Z_k(s)$.

\textbf{1. Tierra virtual} (R2):
\begin{equation}\label{eq:gen-vm}
  V_-(s)=V_+(s)=0 .
\end{equation}

\textbf{2. Ohm en $Z_1$ y Kirchhoff en el nudo inversor} ($I_-=0$, R1):
\begin{equation}\label{eq:gen-I}
  I_2(s)=I_1(s)=\frac{\Vin(s)}{Z_1(s)} .
\end{equation}

\textbf{3. Ohm en $Z_2$:}
\begin{equation}\label{eq:gen-vout}
  \Vout(s)=V_-(s)-Z_2(s)\,I_2(s)=-\frac{Z_2(s)}{Z_1(s)}\,\Vin(s) .
\end{equation}

\textbf{Resultado.} Con $s=\jj\omega$ se obtiene la ganancia en régimen sinusoidal (\cref{prop:sinusoide}):
\begin{equation}\label{eq:gen-H}
  \resultado{H(s)=\frac{\Vout(s)}{\Vin(s)}=-\frac{Z_2(s)}{Z_1(s)}},\qquad
  A_V(\jj\omega)=-\frac{Z_2(\jj\omega)}{Z_1(\jj\omega)} .
\end{equation}
Si $Z_1=R_1$ y $Z_2=R_2$ se recupera \cref{eq:inv-Av}, independiente de $\omega$; si alguna impedancia
contiene un condensador, $Z_C=1/(\jj\omega C)$ introduce la dependencia en frecuencia.
\end{resolucion}

% -------------------------------------------------------------------------
\subsection{Caso $Z_1=R$, $Z_2=R_F\parallel C$: filtro paso bajo}\label{sec:pl-pasobajo}
% -------------------------------------------------------------------------
\begin{resolucion}
\textbf{1. Impedancias} (paralelo, \cref{eq:RC-paralelo}):
\begin{equation}\label{eq:pb-Z}
  Z_1=R,\qquad
  Z_2(s)=\frac{R_F\cdot\frac{1}{sC}}{R_F+\frac{1}{sC}}=\frac{R_F}{1+sR_FC} .
\end{equation}

\textbf{2. Función de transferencia}, sustituyendo \cref{eq:pb-Z} en \cref{eq:gen-H}:
\begin{equation}\label{eq:pb-H}
  \resultado{H(s)=-\frac{R_F}{R}\,\frac{1}{1+sR_FC}}
\end{equation}
Es una red de constante temporal única paso bajo, \cref{eq:STC}, con
\begin{equation}\label{eq:pb-K}
  K=-\frac{R_F}{R},\qquad \tau=R_FC,\qquad f_c=\frac{1}{2\pi R_FC} .
\end{equation}

\textbf{3. Respuesta en frecuencia} ($s=\jj2\pi f$, \cref{eq:STC-modulo}):
\begin{equation}\label{eq:pb-modulo}
  \abs{A_V(f)}=\frac{R_F/R}{\sqrt{1+(f/f_c)^{2}}},\qquad
  \arg A_V(f)=180^{\circ}-\arctg\frac{f}{f_c} .
\end{equation}

\textbf{4. Límites.}
\begin{align}
  f\ll f_c:&\quad \abs{A_V}\simeq\frac{R_F}{R}
     &&\text{($C$ abierto: constante)},\label{eq:pb-baja}\\
  f\gg f_c:&\quad \abs{A_V}\simeq\frac{R_F}{R}\,\frac{f_c}{f}=\frac{1}{2\pi fRC}
     &&\text{($C$ cortocircuita a $R_F$: decrece, $-20\dB$/década)},\label{eq:pb-alta}
\end{align}
y en $f=f_c$, $\abs{A_V}=(R_F/R)/\sqrt2$ ($-3\dB$). Es la dependencia de la figura~4b del guion
(\cref{fig:bode-filtro}): el circuito es un \emph{filtro paso bajo de primer orden}.
\end{resolucion}

% Figura: Bode del filtro paso bajo activo con R=1k, RF=100k, C=1,59 nF (fc=1 kHz).
\begin{figure}[ht]
\centering
\begin{tikzpicture}
  \begin{semilogxaxis}[tema, name=mod, width=0.86\textwidth, height=5.6cm,
      xmin=10, xmax=1e6, ymin=-22, ymax=48, ytick={-20,0,20,37,40},
      xlabel={$f$ (Hz)}, ylabel={$\abs{A_V}$ (dB)}, log basis x=10,
      title style={font=\small}, title={(a) Módulo}, legend pos=south west]
    \addplot[azulacero, line width=1.1pt, domain=10:1e6, samples=300]
      {40-10*log10(1+(x/1001)^2)};
    \addlegendentry{{set 1 ($R_F=100\,\mathrm{k}\Omega$)}}
    \addplot[rojoteo, line width=0.9pt, densely dashed, domain=10:1e6, samples=2]
      {20*log10(100097/x)};
    \addlegendentry{{set 2 ($R_F\to\infty$, integrador)}}
    \addplot[asintota, domain=10:1001, samples=2] {40};
    \addplot[punteado] coordinates {(1001,-22) (1001,37)};
    \addplot[punteado] coordinates {(100097,-22) (100097,0)};
    \addplot[only marks, mark=*, mark size=1.6pt, azulacero] coordinates {(1001,37)};
    \addplot[only marks, mark=*, mark size=1.6pt, rojoteo] coordinates {(100097,0)};
    \node[font=\scriptsize, anchor=south west] at (axis cs:1150,37.5) {$-3\dB$ en $f_c\approx1\,\mathrm{kHz}$};
    \node[font=\scriptsize, anchor=south west] at (axis cs:1.1e5,0.5) {$f_1\approx100\,\mathrm{kHz}$};
    \node[font=\scriptsize, anchor=south west, text=grispizarra, rotate=-24] at (axis cs:6000,22) {$-20\dB$/década};
  \end{semilogxaxis}
  \begin{semilogxaxis}[tema, at={(mod.south)}, anchor=north, yshift=-1.4cm,
      width=0.86\textwidth, height=4.2cm,
      xmin=10, xmax=1e6, ymin=80, ymax=190, ytick={90,135,180},
      xlabel={$f$ (Hz)}, ylabel={$\arg A_V$ (grados)}, log basis x=10,
      title style={font=\small}, title={(b) Fase}]
    \addplot[azulacero, line width=1.1pt, domain=10:1e6, samples=300] {180-atan(x/1001)};
    \addplot[rojoteo, line width=0.9pt, densely dashed, domain=10:1e6, samples=2] {90};
    \addplot[only marks, mark=*, mark size=1.6pt, azulacero] coordinates {(1001,135)};
    \node[font=\scriptsize, anchor=west] at (axis cs:1200,137) {$135^{\circ}$};
  \end{semilogxaxis}
\end{tikzpicture}
\caption{Diagrama de Bode del filtro de la \figpanel{fig:generalizado}{b} con los valores del set~1,
según \cref{eq:pb-modulo}: meseta de $40\dB$ (ganancia $100$), corte a $-3\dB$ en $f_c\approx1\,\mathrm{kHz}$ y
caída de $-20\dB$/década. En trazo discontinuo, el integrador ideal del set~2, que coincide con la
asíntota de alta frecuencia y corta $0\dB$ en $f_1=1/(2\pi RC)\approx100\,\mathrm{kHz}$. La fase pasa de
$180^{\circ}$ (inversión) a $90^{\circ}$ (integración), con $135^{\circ}$ en $f_c$.}
\label{fig:bode-filtro}
\end{figure}


% =========================================================================
\section{Resumen de resultados del PreLab}\label{sec:resumen}
% =========================================================================
La \cref{tab:resumen} reúne los resultados que deben figurar, enmarcados, en el PreLab manuscrito.

\begin{table}[ht]
\centering\small
\renewcommand{\arraystretch}{1.45}
\caption{Resultados de los tres PreLab con el modelo ideal del AO.}
\label{tab:resumen}
\begin{tabular}{@{}lccc@{}}
\toprule
 & Inversor & No inversor & Inversor generalizado\\
\midrule
Ganancia & $A_V=-\dfrac{R_2}{R_1}$ \ \eqref{eq:inv-Av} & $A_V=1+\dfrac{R_2}{R_1}$ \ \eqref{eq:ni-Av}
         & $H(s)=-\dfrac{Z_2}{Z_1}$ \ \eqref{eq:gen-H}\\
Diseño ($R_1=1\,\mathrm{k}\Omega$) & $R_2=10\,\mathrm{k}\Omega$ \ \eqref{eq:inv-R2} & $R_2=10\,\mathrm{k}\Omega$ \ \eqref{eq:ni-R2}
       & $Z_1=R$, $Z_2=R_F\parallel C$\\
Resultado & $A_V=-10$ & $A_V=+11$ & $H(s)=-\dfrac{R_F}{R}\,\dfrac{1}{1+sR_FC}$ \ \eqref{eq:pb-H}\\
\bottomrule
\end{tabular}
\end{table}

Los tres circuitos se resuelven con las mismas dos reglas del AO ideal: corrientes de entrada nulas
y cortocircuito virtual. La ganancia la fija la red pasiva de realimentación, no el AO. Si la red
contiene un condensador, la misma deducción, hecha con impedancias en el dominio de Laplace, da un
filtro paso bajo: ganancia $R_F/R$ por debajo de $f_c=1/(2\pi R_FC)$ y caída de $-20\dB$/década por encima.


% =========================================================================
%                              BIBLIOGRAFÍA
% =========================================================================
\printbibliography[title={Referencias}]

\end{document}
